\documentclass[11pt]{article}
\usepackage[margin=1in]{geometry}
\usepackage{amsmath,amssymb}
\usepackage[hidelinks]{hyperref}

\begin{document}

\section*{Human Review: radial-twist sum-coordinate counterexample}

\noindent Original automatic proof: \texttt{solution\_packet.pdf}

\noindent Checked date: 10 June 2026

\noindent Status: Manually checked.

\noindent Verdict: Appears correct.

\section*{Human Comments}

The solution appears to be correct. The central idea is to encode the
information of a vector \(x\) into two coordinates, with the split depending
on the radius \(\|x\|\). Concretely, the construction has the form
\[
  F(x)=\bigl(\cos(\theta(\|x\|))x,\ \sin(\theta(\|x\|))x\bigr).
\]
Thus the full map \(F=(F_1,F_2)\) still remembers \(x\) in the
two-coordinate sum, but each individual coordinate alternately loses all
information on carefully chosen radial shells.

The alternation is governed by the norm of \(x\). On one sequence of radii,
\(\theta(\|x\|)=0\), so the second coordinate vanishes:
\[
  F_2(x)=0.
\]
On another sequence of radii, \(\theta(\|x\|)=\pi/2\), so the first
coordinate vanishes:
\[
  F_1(x)=0.
\]
The full map remains a coarse Lipschitz embedding because the information is
not destroyed; it is merely moved between the two coordinates as the radius
changes.

\section*{Strength of the Result}

The packet proves something stronger than what is needed for the stated
problem. It does not merely show that there is no infinite-dimensional
subspace on which one coordinate becomes a coarse Lipschitz embedding.
Rather, for every nonzero linear subspace \(X_0\), neither
\[
  F_1|_{X_0}\quad\text{nor}\quad F_2|_{X_0}
\]
is a coarse embedding.

Indeed, after choosing any unit vector \(u\in X_0\), the points \(R_k u\) and
\(-R_k u\) become arbitrarily far apart in \(X_0\), while one of the
coordinate maps sends both to the same point. The same argument applies to
the other coordinate along the alternating sequence of radii. Thus the
failure already occurs on every one-dimensional subspace, not merely on
infinite-dimensional ones.

\section*{Short Version}

The proof works by a radial twist. The vector \(x\) is split between two
coordinates according to \(\|x\|\), and the active coordinate alternates at
larger and larger scales. The pair \(F=(F_1,F_2)\) is still a coarse
Lipschitz embedding, but each separate coordinate collapses arbitrarily
distant pairs in every nonzero subspace.

\section*{Literature and Follow-Up}

We did not find this result stated in the literature during the present
review. Since the surrounding problem is specific and the construction is
quite natural once seen, it may be good to reach out to the author to verify
whether the example is already known and to get his take on its significance.

\section*{Review Outcome}

Archive this packet as an apparently correct negative solution to the
coarse-sum coordinate problem, with the additional note that the construction
is stronger than the stated obstruction: no coordinate map works even after
restriction to a nonzero finite-dimensional subspace.

\end{document}
